% !TeX root = ../navier-stokes-manuscript.tex
\subsection{The heat clock}\label{sub:BodyHeatClock}

The scaling calculation tells us that a smaller vortex can remain just as critical.
It does not tell us how much time the fluid has to make the next smaller state.

A structure of width \(\ell\) is carried by frequencies of order \(|\xi|\sim\ell^{-1}\).
Under the linear heat flow, those frequencies decay according to
\[
  \widehat v(\xi,t+\delta t)
  =e^{-\nu|\xi|^2\delta t}\widehat v(\xi,t).
\]
Viscosity makes an order-one change when \(\nu|\xi|^2\delta t\) is of order one.
At scale \(\ell\), this gives
\begin{equation}\label{eq:BodyHeatComparisonTime}
  \delta t_{\mathrm{heat}}(\ell)
  \asymp\frac{\ell^2}{\nu}.
\end{equation}
This is the heat clock.
When the vortex narrows from \(\ell\) to \(\ell/\lambda\), its viscous response time shrinks by \(\lambda^{-2}\).

Now repeat the contraction at a fixed ratio:
\[
  \ell_k=\rho^k\ell_0,
  \qquad
  0<\rho<1.
\]
The heat time attached to the \(k\)-th scale is
\[
  \delta t_k^{\mathrm{heat}}
  \asymp
  \frac{\ell_0^2}{\nu}\rho^{2k},
\]
and the whole sequence satisfies
\[
  \sum_{k=0}^{\infty}\delta t_k^{\mathrm{heat}}
  \asymp
  \frac{\ell_0^2}{\nu}
  \sum_{k=0}^{\infty}\rho^{2k}
  =
  \frac{\ell_0^2}{\nu(1-\rho^2)}
  <\infty.
\]
Every new scale comes with less time than the one before it, and all of those comparison times fit inside one bounded interval.
That is the Zeno possibility.

The geometric series by itself is only a clock.
The fluid still has to use those shrinking intervals to reach more dangerous states.
An exact rescaled copy leaves \(H\) unchanged, so contraction alone cannot produce critical-height blow-up.
A realized cascade needs both parts at once: the active length must fall while the critical height rises.

Choose \(H_0>H(0)\) and set
\[
  L_k:=2^kH_0.
\]
Whenever the level is reached, let
\[
  t_k:=\inf\{t>0:H(t)=L_k\}.
\]
If arbitrarily high levels are reached, continuity orders the first passages:
\[
  t_0<t_1<t_2<\cdots.
\]
Suppose the \(k\)-th passage occurs at the active length \(\ell_k=\rho^k\ell_0\), and suppose the next passage is completed within a fixed multiple of the heat time:
\[
  0<t_{k+1}-t_k
  \leq
  C\frac{\ell_k^2}{\nu}
  =
  C\frac{\ell_0^2}{\nu}\rho^{2k}.
\]
Then
\[
  \sum_{k=0}^{\infty}(t_{k+1}-t_k)<\infty,
  \qquad
  t_k\uparrow\tau
\]
for some \(\tau<\infty\).
The fluid is now crossing higher critical levels while the active scale and the interval available for each crossing collapse together.
This is the Zeno cascade attached to critical height.

The same rescaling sharpens spatial derivatives:
\[
  \nabla^m u_\lambda(x,t)
  =\lambda^{m+1}(\nabla^m u)(\lambda x,\lambda^2t),
  \qquad
  m\in\Nzero.
\]
At the \(k\)-th active scale, \(\lambda_k:=\ell_0/\ell_k=\rho^{-k}\), so the first gradient carries the factor \(\lambda_k^2\).
If the normalized gradients of the realized states remain nonzero, the physical gradients diverge as \(k\to\infty\).
The scale is collapsing, the crossings are accelerating, and the gradients are sharpening without bound.
Their limit is the singularity described by the critical-height vortex.

Between one crossing and the next, the velocity still has to make a real change.
Once those intervals collapse, that change becomes a demand on acceleration, jerk, and every later response.
